% !TEX program = xelatex
% 공통수학2 · Ⅲ. 함수와 그래프 · 1. 함수 · 01 함수 · 2/2차시
\documentclass[11pt,a4paper]{article}
\usepackage{kotex}
\usepackage{iftex}
\ifPDFTeX\else
  \setmainhangulfont{Noto Serif CJK KR}
  \setsanshangulfont{Noto Sans CJK KR}
\fi
\usepackage[margin=2.1cm]{geometry}
\usepackage{parskip}
\usepackage{enumitem}
\usepackage{array}
\usepackage{tabularx}
\usepackage{booktabs}
\usepackage{xcolor}
\usepackage{amssymb}
\usepackage{tikz}
\usepackage[most]{tcolorbox}
\usepackage{fancyhdr}
\usepackage{titlesec}

\definecolor{goldc}{HTML}{B07C22}
\definecolor{goldstrong}{HTML}{8A5F16}
\definecolor{indigoc}{HTML}{2B3B57}
\definecolor{indigosoft}{HTML}{6E82A6}
\definecolor{linec}{HTML}{D6DACB}
\definecolor{surface2}{HTML}{E4E8DC}
\definecolor{notebg}{HTML}{FBF3E3}
\definecolor{inksoft}{HTML}{52604F}

\titleformat{\section}{\normalfont\sffamily\bfseries\large\color{indigoc}}{\thesection}{0.6em}{}
\titlespacing{\section}{0pt}{16pt}{8pt}
\newtcolorbox{ideabox}{enhanced, breakable, colback=white, colframe=goldc, boxrule=0.5pt, leftrule=3pt, arc=2pt, left=10pt, right=10pt, top=8pt, bottom=8pt}
\newtcolorbox{calloutbox}{enhanced, breakable, colback=white, colframe=indigosoft, boxrule=0.5pt, leftrule=3pt, arc=2pt, left=10pt, right=10pt, top=7pt, bottom=7pt}
\newtcolorbox{notebox}{enhanced, breakable, colback=notebg, colframe=goldc, boxrule=0.5pt, leftrule=3pt, arc=2pt, left=10pt, right=10pt, top=7pt, bottom=7pt}
\newtcolorbox{answerbox}{enhanced, breakable, colback=white, colframe=linec, boxrule=0.5pt, arc=2pt, left=10pt, right=10pt, top=7pt, bottom=7pt}
\newcommand{\lessonbanner}[3]{%
  \begin{tcolorbox}[enhanced, colback=indigoc, colframe=indigoc, boxrule=0pt, arc=0pt,
    left=14pt, right=14pt, top=14pt, bottom=10pt, borderline south={2.4pt}{0pt}{goldc}, fontupper=\color{white}]
    {\sffamily\footnotesize\color{white!85!black} #1}\\[4pt]{\sffamily\bfseries\Large #2}\\[3pt]{\sffamily\small\color{white!88!black} #3}
  \end{tcolorbox}}
\pagestyle{fancy}\fancyhf{}\renewcommand{\headrulewidth}{0pt}
\fancyfoot[L]{\footnotesize\color{inksoft} 공통수학2 · 2026학년도 2학기 · 지평선고등학교}
\fancyfoot[R]{\footnotesize\color{inksoft} 교사용 지도안 -- 배포 금지}

\newcommand{\mapfour}[1]{%
  \begin{tikzpicture}[scale=0.58, baseline]
    \draw[thick, indigosoft] (0,0) ellipse (0.9 and 1.6);
    \draw[thick, indigosoft] (2.6,0) ellipse (0.9 and 1.6);
    \node[circle, fill=indigoc, inner sep=1.4pt, label=left:{\footnotesize $1$}] (x1) at (0,1) {};
    \node[circle, fill=indigoc, inner sep=1.4pt, label=left:{\footnotesize $2$}] (x2) at (0,0) {};
    \node[circle, fill=indigoc, inner sep=1.4pt, label=left:{\footnotesize $3$}] (x3) at (0,-1) {};
    \node[circle, fill=goldstrong, inner sep=1.4pt, label=right:{\footnotesize $1$}] (y1) at (2.6,1) {};
    \node[circle, fill=goldstrong, inner sep=1.4pt, label=right:{\footnotesize $3$}] (y2) at (2.6,0) {};
    \node[circle, fill=goldstrong, inner sep=1.4pt, label=right:{\footnotesize $5$}] (y3) at (2.6,-1) {};
    \node[above] at (0,1.75) {\footnotesize $X$};
    \node[above] at (2.6,1.75) {\footnotesize $Y$};
    #1
  \end{tikzpicture}}

\newcommand{\mapthreefour}[1]{%
  \begin{tikzpicture}[scale=0.58, baseline]
    \draw[thick, indigosoft] (0,0) ellipse (0.9 and 1.9);
    \draw[thick, indigosoft] (2.8,0) ellipse (0.9 and 1.9);
    \node[circle, fill=indigoc, inner sep=1.4pt, label=left:{\footnotesize $1$}] (x1) at (0,1.2) {};
    \node[circle, fill=indigoc, inner sep=1.4pt, label=left:{\footnotesize $2$}] (x2) at (0,0) {};
    \node[circle, fill=indigoc, inner sep=1.4pt, label=left:{\footnotesize $3$}] (x3) at (0,-1.2) {};
    \node[circle, fill=goldstrong, inner sep=1.4pt, label=right:{\footnotesize $1$}] (y1) at (2.8,1.8) {};
    \node[circle, fill=goldstrong, inner sep=1.4pt, label=right:{\footnotesize $3$}] (y2) at (2.8,0.6) {};
    \node[circle, fill=goldstrong, inner sep=1.4pt, label=right:{\footnotesize $5$}] (y3) at (2.8,-0.6) {};
    \node[circle, fill=goldstrong, inner sep=1.4pt, label=right:{\footnotesize $7$}] (y4) at (2.8,-1.8) {};
    \node[above] at (0,2.1) {\footnotesize $X$};
    \node[above] at (2.8,2.1) {\footnotesize $Y$};
    #1
  \end{tikzpicture}}

\newcommand{\gridaxes}{%
  \draw[step=1cm, linec, very thin] (-2.2,-2.2) grid (2.2,2.2);
  \draw[->, thick, indigosoft] (-2.2,0) -- (2.5,0) node[right, font=\small]{$x$};
  \draw[->, thick, indigosoft] (0,-2.2) -- (0,2.5) node[above, font=\small]{$y$};}

\begin{document}
\lessonbanner{공통수학2 · Ⅲ. 함수와 그래프 · 1. 함수}{01 함수 --- 2/2차시}{함수의 그래프와 여러 가지 함수 · 교사용 지도안 (2026학년도 2학기)}
\vspace{10pt}

\noindent\begin{tabularx}{\linewidth}{@{}>{\bfseries\color{indigoc}}p{2.6cm}X@{}}
\toprule
대상 & 고1 · 2개 반 · 반당 13명 \\
차시 & 50분 (도입 10 · 전개 30 · 정리 10) \\
성취기준 & 함수의 그래프를 이해하고, 여러 가지 함수(일대일함수, 일대일대응, 항등함수, 상수함수)를 구별할 수 있다. \\
선행 학습 & 33차시 --- 함수의 뜻, 정의역·공역·치역 \\
\bottomrule
\end{tabularx}

\section{학습 목표 \& Big Idea}
\begin{ideabox}
\textbf{\color{goldstrong}영속적 이해 (Big Idea)}\\
함수의 그래프는 대응이라는 추상적 규칙을 좌표평면 위의 `모양'으로 번역한 것이다. 그리고 대응이 갖는 특별한 성질(하나씩만 대응하는가, 공역이 남김없이 채워지는가)에 따라 함수는 서로 다른 종류로 분류된다.
\vspace{4pt}
\small\color{inksoft} 핵심개념: \textbf{그래프의 판별} \quad\textbullet\quad 핵심개념: \textbf{일대일성과 전사성} \quad\textbullet\quad 일반화: \textbf{대응의 성질이 함수의 종류를 결정한다}
\end{ideabox}
\begin{calloutbox}
\textbf{차시 학습 목표}
\begin{itemize}[leftmargin=1.4em, itemsep=2pt, topsep=4pt]
  \item 함수의 그래프의 뜻을 알고, 수직선 판정법으로 어떤 그래프가 함수의 그래프인지 판별할 수 있다.
  \item 일대일함수, 일대일대응, 항등함수, 상수함수를 구별하고 예를 들 수 있다.
\end{itemize}
\end{calloutbox}

\section{질문의 연결}
\begin{tabularx}{\linewidth}{@{}>{\bfseries\color{indigoc}\small}p{2.4cm}X@{}}
사실적 질문 & 좌표평면 위의 어떤 도형이 함수의 그래프가 되려면 어떤 조건을 만족해야 할까? \\[6pt]
개념적 질문 & 일대일함수와 일대일대응은 무엇이 같고 무엇이 다를까? \\[6pt]
논쟁적 질문 & 우리 반에서 `번호 $\to$ 그 번호 학생의 키'라는 대응은 일대일함수일까? 그렇게 되려면 어떤 조건이 필요할까? \\
\end{tabularx}

\section{수업 흐름}
\begin{tabularx}{\linewidth}{@{}>{\centering\arraybackslash}p{1.6cm}X>{\raggedright\arraybackslash\small}p{3.1cm}@{}}
\toprule
\textbf{단계} & \textbf{활동 \& 발문} & \textbf{ATL / VTR} \\
\midrule
\textcolor{goldstrong}{\textbf{도입}}\newline\footnotesize 10분 &
\textbf{마음공부 (2분)} --- 유무념 대조.\newline
\textbf{생각 열기 (8분)} --- 지난 시간 배운 $f:X\to Y$를 좌표평면 위의 점 $(x,f(x))$로 하나씩 찍어보며 `그래프란 대응을 그림으로 옮긴 것'임을 확인. &
ATL: 사고(표상 전환)\newline VTR: See-Think-Wonder \\
\addlinespace
\textcolor{goldstrong}{\textbf{전개}}\newline\footnotesize 30분 &
\textbf{그래프의 판별 (10분)} --- 함수의 그래프 정의 $\{(x,f(x))\mid x\in X\}$ 소개 $\to$ 포물선과 원 그래프 제시, 수직선(세로선) 판정법 유도(정의역의 각 $x$에 $y$값이 오직 하나여야 함) $\to$ 문제 1.\newline
\textbf{여러 가지 함수 (12분)} --- 일대일함수(공역의 원소가 두 번 이상 쓰이지 않음) $\to$ 일대일대응(일대일함수이면서 치역=공역) $\to$ 항등함수 $f(x)=x$ $\to$ 상수함수 $f(x)=c$를 화살표 그림과 식으로 각각 제시 $\to$ 문제 2, 문제 3.\newline
\textbf{수평선 판정법 (8분)} --- 일대일함수의 그래프는 임의의 수평선과 많아야 한 점에서 만남을 소개 $\to$ 문제 4. &
ATT: 촉진자 역할\newline ATL: 분류·비교 \\
\addlinespace
\textcolor{goldstrong}{\textbf{정리}}\newline\footnotesize 10분 &
\textbf{개념 연결 질문 (5분)} --- ``일대일함수인데 일대일대응이 아닌 예를 하나 만들어 설명해 보자.''\newline
\textbf{성찰 (5분)} --- 감사 일기 1문장 + 다음 차시 예고(합성함수). &
마음공부 · 감사 일기 \\
\bottomrule
\end{tabularx}

\begin{calloutbox}
\textbf{지도상의 유의점} --- 일대일함수와 일대일대응을 혼동하기 쉽다. 일대일함수는 `공역의 원소가 두 번 이상 대응되지 않음'(유일성 + 각 상이 서로 다름)만 요구하고, 일대일대응은 여기에 `치역과 공역이 일치함'(전사성)까지 추가로 요구한다는 점을 화살표 그림으로 명확히 대조한다.
\end{calloutbox}

\section{형성평가 \& 답안}
\begin{minipage}[t]{0.48\linewidth}
\begin{answerbox}
\textbf{문제 1} \quad 다음 두 그래프 중 함수의 그래프인 것을 고르고 이유를 설명하시오.

\vspace{4pt}
\begin{tikzpicture}[scale=0.5]
  \gridaxes
  \draw[thick, goldc, domain=-1.85:1.85, smooth, variable=\t] plot({\t},{0.5*\t*\t-1});
  \draw[dashed, inksoft] (1,-2.2) -- (1,2.2);
  \node[below] at (0,-2.6) {\footnotesize ① $y=\frac12x^2-1$};
\end{tikzpicture}\hfill
\begin{tikzpicture}[scale=0.5]
  \gridaxes
  \draw[thick, goldc] (0,0) circle (1.5);
  \draw[dashed, inksoft] (0.8,-2.2) -- (0.8,2.2);
  \node[below] at (0,-2.6) {\footnotesize ② $x^2+y^2=2.25$};
\end{tikzpicture}

\vspace{4pt}
\textbf{\color{goldstrong}답} \; ①. 어떤 세로 직선(수직선)을 그어도 그래프와 항상 한 점에서만 만난다. ②는 원의 위쪽·아래쪽 두 점에서 만나는 세로 직선이 있으므로 하나의 $x$에 $y$값이 두 개 대응 --- 함수 아님.
\end{answerbox}
\end{minipage}\hfill
\begin{minipage}[t]{0.48\linewidth}
\begin{answerbox}
\textbf{문제 2} \quad $X=\{1,2,3\}$, $Y=\{1,3,5,7\}$, $f(x)=2x-1$의 대응을 그림으로 나타내고, 일대일함수인지 일대일대응인지 판별하시오.

\vspace{4pt}
\centering
\mapthreefour{
  \draw[->, thick] (x1) -- (y1);
  \draw[->, thick] (x2) -- (y2);
  \draw[->, thick] (x3) -- (y3);
}

\vspace{4pt}
\raggedright\textbf{\color{goldstrong}답} \; 일대일함수(서로 다른 원소의 함숫값이 모두 다름). 그러나 $7\in Y$가 대응되지 않아 치역 $\{1,3,5\}\neq$ 공역 $\{1,3,5,7\}$이므로 일대일대응은 아니다.
\end{answerbox}
\end{minipage}

\vspace{6pt}
\begin{minipage}[t]{0.48\linewidth}
\begin{answerbox}
\textbf{문제 3} \quad $X=\{1,2,3\}$일 때, 항등함수 $f(x)=x$와 상수함수 $g(x)=2$의 함숫값을 각각 구하시오.\\[4pt]
\textbf{\color{goldstrong}답} \; $f(1)=1,\,f(2)=2,\,f(3)=3$ (자기 자신) \quad $g(1)=g(2)=g(3)=2$ (항상 같은 값)
\end{answerbox}
\end{minipage}\hfill
\begin{minipage}[t]{0.48\linewidth}
\begin{answerbox}
\textbf{문제 4} \quad $y=\frac12x^2-1$의 그래프는 일대일함수의 그래프인가? 수평선 판정법으로 설명하시오.\\[4pt]
\textbf{\color{goldstrong}답} \; 아니다. $y=1$과 같은 수평선을 그으면 그래프와 두 점(대칭인 두 $x$값)에서 만나므로 서로 다른 $x$에 같은 함숫값이 대응 --- 일대일함수 아님.
\end{answerbox}
\end{minipage}

\section{성취수준 (이 차시 범위)}
\begin{tabularx}{\linewidth}{@{}>{\bfseries\color{goldstrong}}p{0.9cm}X@{}}
\toprule
A & 수직선·수평선 판정법의 원리를 설명하고, 여러 가지 함수를 근거를 들어 정확히 구별하며 반례를 스스로 구성할 수 있다. \\
B & 함수의 그래프를 판별하고, 일대일함수·일대일대응·항등함수·상수함수를 구별할 수 있다. \\
C & 안내에 따라 함수의 그래프를 판별하고, 여러 가지 함수의 뜻을 안다. \\
D & 항등함수와 상수함수의 함숫값을 구할 수 있다. \\
E & 함수의 그래프가 무엇인지 안다. \\
\bottomrule
\end{tabularx}

\section{다음 차시}
\begin{calloutbox}
\textbf{02 합성함수 --- 1/1차시.} 두 함수를 이어 붙여 새로운 함수를 만드는 합성함수를 다룬다.
\end{calloutbox}
\end{document}
